\chapter*{Foreword and acknowledgments}


本书是一份旨在用于操作系统课程的草稿文本。
它通过研究一个名为 xv6 的示例内核，来解释操作系统的主要概念。
Xv6 是基于 Dennis Ritchie 和 Ken Thompson 的 Unix Version 6 (v6)~\cite{unix} 建模的。
Xv6 在结构和风格上大致遵循了 v6，但它是用 ANSI C~\cite{kernighan} 实现的，并针对多核 RISC-V~\cite{riscv}。
本书应当与 xv6 的源代码结合阅读，这种方法受到了 John Lions' Commentary on UNIX 6th Edition~\cite{lions} 的启发；
本文中包含指向位于 \url{https://github.com/mit-pdos/xv6-riscv} 的源代码的超链接。请访问 \url{https://pdos.csail.mit.edu/6.1810} 以获取有关 v6 和 xv6 的在线资源的更多指针，其中包括几个使用 xv6 的实验作业。
我们已经在 MIT 的操作系统课程 6.828 和 6.1810 中使用了本书。
我们感谢这些课程的教职人员、助教和学生，他们都直接或间接地对 xv6 做出了贡献。
特别地，我们要感谢 Adam Belay、Austin Clements 和 Nickolai Zeldovich。
最后，我们要感谢那些通过电子邮件向我们指出书中错误或提出改进建议的人：Abutalib Aghayev、Sebastian Boehm、brandb97、Anton Burtsev、Raphael Carvalho、Tej Chajed、Brendan Davidson、Rasit Eskicioglu、Color Fuzzy、Wojciech Gac、Giuseppe、Tao Guo、Haibo Hao、Naoki Hayama、Chris Henderson、Robert Hilderman、Eden Hochbaum、Wolfgang Keller、Paweł Kraszewski
